\chapter{DOS Tutor Evidence Record}

For each DOS-era program documented in Part~VI, record the following.

\begin{itemize}
\item Title, publisher, edition, and platform.
\item Screenshots or reconstructed layouts of the menu, a lesson screen, a game screen, and a results screen.
\item Manual text describing lessons, thresholds, and scoring.
\item The complete exercise corpus for at least the first ten lessons, transcribed.
\item Scoring rules and any advancement thresholds.
\item A recorded emulated session with the ledger schema of Chapter~27 applied to it.
\item The source of each item and the date it was obtained.
\end{itemize}

\section*{Records so far}

\begin{itemize}
\item \textbf{Typing Tutor III with Letter Invaders} (Kriya Systems, 1984; Simon \& Schuster). Have: title and menu screen (S1), and catalog captions listing the screens of the IBM PC version (users, practice, status reports, graphs, game overview and summary). Need: the screenshots themselves, the manual, exercise corpus, scoring rules, and the level structure behind the S5 key-order strip.
\item \textbf{Typing Tutor IV with Letter Invaders} (Kriya Systems, 1983 and 1987 on the title screen; Simon \& Schuster Software). Have: title screen (S2) and a lesson screen (S3), both from the v1.0 (1987) single-disk image described on the Internet Archive; library and retail descriptions of its adaptive behavior. Need: an emulated session recorded with the ledger schema, the manual, gauge labels, metronome settings, the Letter Invaders screen, and the lesson-selection rule.
\item \textbf{Typing Tutor III, further screens}: S5, the letters game with its key-order strip, is from this program (1984). Need: the level structure behind the strip.
\item \textbf{Unidentified falling-words game} (S4). Need: identification.
\end{itemize}

\chapter{Scoring Rules}

\section{Speed}

Gross and net speeds are defined in Chapter~17 with the five-character word. Timed passages in the workbook print a cumulative word count at the right of each line, so the word count at the end of a timing is read from the last completed line plus the count for the partial line, and speed is that count divided by the five minutes of the timing.

\section{Workbook Drill Scale}

On drill lines the errors in each set of lines are marked and counted, and the count is read against a scale.

\begin{center}
\begin{tabular}{@{}ll@{}}
\toprule
Errors per set of lines & Reading \\
\midrule
0 to 3 & excellent \\
4 to 7 & satisfactory \\
8 or more & extra practice \\
\bottomrule
\end{tabular}
\end{center}

\section{Timed Writings}

Timed writings are five minutes, practiced once, then timed and repeated twice, scored by the supervisor. The thresholds applied at Progress Tests are not known from the pages available. \needs{Progress Test page and Lesson Completion Sheet}

\chapter{Exercise Transcriptions}

Each group is printed as three identical lines of the given unit, cycled to the line width and ending mid-unit. Units are transcribed once; line endings are noted where they differ. Lines that could be read only as show-through are omitted.

\section{Academy Pages, Exercise 1 (p.\ 11), keys A S D F ; L K J}

\begin{Verbatim}
G1  asdf ;lkj                     (ends: asd)
G2  asdfa ;lkj;                   (ends: asd)
G3  dafs k;jl  dafs k;jk  dafs k;jl  dafs k;jl ...   (ends: daf)
G4  ass lad add alf               (last unit: ass lad add ass)
G5  flak alas fall asks           (ends: ask)
\end{Verbatim}

\section{Academy Pages, Exercise 2 (p.\ 12), new keys G H}

\begin{Verbatim}
G1  asdfg ;lkjh                   (ends: asd)
G2  asdfgfa ;lkjhj;               (last unit: ;lk asd)
G3  gaff gash lags hall           (ends: lag)
G4  glad asks hash gala gald asks hash gala   (ends: ask)
G5  salad glass shall flash jaffa slash       (ends: lad)
G6  a sad lad asks dad;           (ends: sad)
\end{Verbatim}

\section{Academy Pages, Exercise 6 (pp.\ 20, 21), new keys W I}

\begin{Verbatim}
G1  asdsws ;lkikl                 (ends: asdsws;)
G2  AsDsWs ;LkIkL                 (ends: As)
G3  wed kid was lie
G4  ail law his owe
G5  were soil will hide           (ends: sod)
G6  heir wear like wore           (ends: her)
G7  wise lied wild wire           (ends: lie)
G8  A wood fire is good; oil is worse for a fire; wood is for fire
G9  We killed wild doe for a while for food; Will would ride well
G10 A fellow said he walked a high hill; he grew riled if we rode;
G11 Ilse who is a redhead is a selfish girl who likes a wild life;
\end{Verbatim}

\section{Academy Pages, Exercise 13 (p.\ 34), new key Z}

\begin{Verbatim}
G1  asaza ;l;p;                   (ends: asa)
G2  aSaZa ;L;P;
G3  zip zoo zig zag
G4  zinc: zulu zero jazz:
G5  quiz doze size hazy:
G6  zeal lazy fizz fuzz:
\end{Verbatim}

\section{Academy Pages, an Exercise Between 6 and 13 (p.\ 31), Groups 7 to 13}

\begin{Verbatim}
G7  prove quote point quart print query      (ends: quest)
G8  quota press equip paste
G9  Peter and Paul ran quickly to the printers for a paper:
G10 King Pepe snapped his fingers as the Pekinese pup passed by;
G11 This lively path is paved and has a border of crimson poppies;
G12 I will quit the room again if he does not answer all queries:
G13 The Queen was presented with two quills at the banquet today;
\end{Verbatim}

\section{Workbook, Lesson 1 (p.\ 1-1)}

\begin{Verbatim}
Speed Drill (each line three times)
 The machine is of the same high quality.
 See that it is cut from the second list.
 The amount of the bank balance is a surprise.
 It will be well to place the paper on record.

Keyboard Drill (each line twice)
 a;sldkfjgh ghfjdksla; a;sldkfjgh ghfjdksla; a;sldkfjgh
 aqa sws ded frf ftf fgf ;p; lol kik juj jyj jhj aza sxs
 dcd fvf fbf ;/; 1.1 k,k jmj jnj
 aqa aqua sws sail ded dead frf rain ftf tall fgf gulf juj
 undo jyj year jhj huge kik kilo lol lock ;p; paws aza zoom
 sxs exited dcd cuts fvf vain fbf bore jnj noon jmj move
 k,k asks 1.1 Rd. ;/; a/c

Alphabet Drill (each set of lines twice)
 a b c d e f g h i j k l m n o p q r s t u v w x y z
 ab cd ef gh ij kl mn op qr st uv wx yz
 abc def ghi jkl mno pqr stu vwx yz
 and but can did eye for get him ice jam kit led me nip oil
 pad quit rob say ton use vat why axe yet zoo
 Only a few quite exceptional people do realize that there is
 very little value in just talking without being able to say
 something worthwhile.

Word Families, Phrases, Common Word Sentence (each line once)
 sell sells sold; like likes liked; make makes made; sew sews
 sewed; turn turns turned; write writes written; act acts acted
 of the/in the/to the/but the/on the/and the/of the/we are/he
 is/they are/I was/you were/it will/it might be/you will find/
 He is planning to travel to the west coast early in October.
 John was able to ask Robert to list all the supplies I need.
 I think those skates will be advertised in the paper Monday.
 There are not that many large parks close to where he lives.
\end{Verbatim}

\chapter{Notation and Procedures}

\section{Notation}

\begin{center}
\begin{tabular}{@{}lp{9.5cm}@{}}
\toprule
Symbol & Meaning \\
\midrule
$\Kset_i$ & set of keys admitted by exercise $i$ \\
$\Newk_i$ & new keys, $\Kset_i\setminus\Kset_{i-1}$ \\
$\Sig(w)$ & set of characters required to type word $w$ \\
$\Vocab_i$ & vocabulary available at exercise $i$ \\
$\Vocab_i^{\mathrm{new}}$ & words in $\Vocab_i$ using at least one new key \\
$G_j$, $E_i$ & group $j$; exercise $i$ \\
$S=(A,V,H,E,T,R)$ & learner state: accuracy, velocity, hesitation, error distribution, transition competence, correction behavior \\
$\tau$ & response time available in a falling-letter game \\
$J$ & advancement judgment \\
\bottomrule
\end{tabular}
\end{center}

\section{Computing the Vocabulary of a Lesson}

The following short script computes $\Vocab_i$ and $\Vocab_i^{\mathrm{new}}$ from a plain word list with one lowercase word per line. The key sets are supplied as strings of characters.

\begin{Verbatim}
import sys

def vocab(words, K):
    return [w for w in words if set(w) <= set(K)]

def vocab_new(words, K, N):
    return [w for w in vocab(words, K) if set(w) & set(N)]

# QWERTY position-to-Dvorak mapping, by physical key
Q = "qwertyuiopasdfghjkl;zxcvbnm,./"
D = "',.pyfgcrlaoeuidhtns;qjkxbmwvz"
to_dvorak = dict(zip(Q, D))

def transpose(text):
    return "".join(to_dvorak.get(c, c) for c in text)

if __name__ == "__main__":
    words = [w.strip() for w in open(sys.argv[1]) if w.strip().isalpha()]
    K1_qwerty = "asdfjkl;"
    K1_dvorak = transpose(K1_qwerty)
    print(len(vocab(words, K1_qwerty)), len(vocab(words, K1_dvorak)))
\end{Verbatim}

The script is a specification of the measurement in Chapter~22. The counts it returns depend on the word list, which must be stated with any reported figure.

\chapter{Glossary}

\begin{description}
\item[Admissible vocabulary.] The words that can be typed using only the keys admitted so far.
\item[Chord.] A set of keys pressed together as one stroke.
\item[Coverage exercise.] A text constructed to contain every letter, such as a full-alphabet sentence.
\item[Ladder.] The sequence of groups in an exercise, from chains of the new keys to sentences.
\item[Legibility.] The property of a tutor whose pedagogical rule can be recovered from its interface.
\item[Position matching.] Transposing a lesson from one layout to another by physical key location.
\item[Sandwich drill.] A drill in which a reached key is placed between two presses of a home-row key.
\item[Stroke.] The unit of action in a chorded system.
\item[Uncorrected error rate.] The fraction of final characters that are wrong.
\item[Corrected error rate.] The fraction of keystrokes that were errors, whether or not repaired.
\end{description}

\chapter{Evidence Register}

\begin{center}
\begin{tabular}{@{}p{4.4cm}p{3.2cm}p{4.2cm}@{}}
\toprule
Source & Status & Note \\
\midrule
Academy pages 1, 11, 12, 20, 21, 31, 34 & photographed, transcribed & exercise number of p.~31 not established; other exercises absent \\
Workbook instruction pages i, ii, iii & photographed & text on the reverse partly legible as show-through \\
Workbook table of contents, Lessons 1 to 20 & photographed & complete \\
Workbook Lesson 1, p.~1-1 & photographed, transcribed & later lessons absent \\
Screenshots S1, S2 & title screens, identified & Typing Tutor III and IV with Letter Invaders \\
PixelatedArcade catalog, Typing Tutor III & captions only, read 20 September 2026 & 1984, IBM PC compatibles, Kriya Systems, Simon \& Schuster \\
Screenshot S3 & lesson screen & Typing Tutor IV v1.0, per the owner \\
Internet Archive item, Typing Tutor IV v1.0 (1987) & catalog description & single 720K disk, DOS, start file tt.exe; date-entry caveat \\
Library and retail descriptions of Typing Tutor IV & secondary & adaptive behavior, registration, editions; to be verified \\
Screenshot S5 & game screen, Typing Tutor III & strip read from the image; meaning inferred \\
Screenshot S4 & game screen & program not identified \\
Handwritten Dvorak retraining notes & photographed & used as evidence in Chapter~24; the meaning of several small annotations was confirmed by the owner \\
Handwritten design notes for Chapter~8 and the QWERTY and Dvorak curriculum sheet & photographed & design, not evidence of any historical tutor \\
\bottomrule
\end{tabular}
\end{center}

\section*{Outstanding Items}

\begin{itemize}
\item Whether the Academy pages and the workbook shared a physical binder, and how a learner moved between them.
\item A Lesson Completion Sheet and a Progress Test page.
\item The reverse sides of Academy pages 20 and 34, and pages for the missing exercises.
\item Later workbook lessons (numbers, symbols, alternate-hand and one-hand words, spelling).
\item A curated, hand-checked word list to replace the web-derived list used in Chapter~22.
\item Identification of screenshot S4.
\item A recorded emulated session of Typing Tutor IV v1.0 and of Typing Tutor III.
\item DOS-era program manuals and further screens, per Appendix~A.
\item Sources for every item flagged in the text.
\end{itemize}

\chapter{Generated Corpora}

These corpora instantiate the ladder template of Chapter~\ref{ch:qwerty} for the design key order. They are generated by a script, not transcribed from any tutor. Words are drawn in frequency order from the google-10000-english list (no-swears version), after removing tokens shorter than three letters (with a short whitelist of common two-letter words) and a blacklist of place names, brands, and fragments. Groups G3 to G5 take words that use at least one new key, ordered by rank, in length bands of three to four, four to five, and six to nine letters. G6 is a string of frequent admissible words and is not a grammatical phrase. The lists are a starting point for editing, and the leftover oddities (for example the ``words'' \texttt{dsl} and \texttt{aka} that survive in the frequency list) are a reminder that a web-derived list has its own admissibility problems. Each exercise prints every group as three lines in the paper tutor; here each group is shown once.

Chains (G1, G2) are sandwich drills using the home key of the same finger column. In the Dvorak corpus the chains are the QWERTY chains transposed by position, and capitals are applied only where the transposed key is a letter.

\section{QWERTY Design Corpus}
\begin{Verbatim}
Exercise 1   new keys: ;adfjkls
  G1  asdf ;lkj
  G2  asdfa ;lkj;
  G3  all add ask fall
  G4  fall falls adds salad
  G5  falls salad
  G6  as all add ask ad

Exercise 2   new keys: gh
  G1  asdfg ;lkjh
  G2  asdfgfa ;lkjhj;
  G3  has had half gas
  G4  shall half flash glass
  G5  shall flash glass flags
  G6  has had half gas hall

Exercise 3   new keys: eu
  G1  ded juj
  G2  dEd jUj
  G3  use see used she
  G4  used full sale sales
  G5  useful league glasses sealed
  G6  us he use see used

Exercise 4   new keys: or
  G1  frf lol
  G2  fRf lOl
  G3  for are free our
  G4  free here also good
  G5  should address offers seller
  G6  of for or are free

Exercise 5   new keys: ty
  G1  ftf jyj
  G2  fTf jYj
  G3  the that you your
  G4  that your other they
  G5  health through results states
  G6  the to that you at

Exercise 6   new keys: iw
  G1  sws kik
  G2  sWs kIk
  G3  this with was will
  G4  this with will site
  G5  software rights details without
  G6  is this with it was

Exercise 7   new keys: mn
  G1  jnj jmj
  G2  jNj jMj
  G3  and not from new
  G4  from more home time
  G5  online system message general
  G6  and in on not from

Exercise 8   new keys: bv
  G1  fvf fbf
  G2  fVf fBf
  G3  have but web view
  G4  have about view been
  G5  business number available review
  G6  by be have but web

Exercise 9   new keys: ,c
  G1  dcd k,k
  G2  dCd k,k
  G3  can back city such
  G4  which click back music
  G5  search contact services service
  G6  can back city such each

Exercise 10   new keys: pq
  G1  aqa ;p;
  G2  aQa ;P;
  G3  page help top post
  G4  page help price post
  G5  people products product policy
  G6  page up help top post

Exercise 11   new keys: .x
  G1  sxs l.l
  G2  sXs l.l
  G3  next text box tax
  G4  next index text linux
  G5  example exchange executive excellent
  G6  next text box tax fax

Exercise 12   new keys: /
  G1  ;/;
  G2  ;/;
  G3  (no listed word contains a new key)
  G4  (no listed word contains a new key)
  G5  (no listed word contains a new key)
  G6  (none)

Exercise 13   new keys: z
  G1  aza
  G2  aZa
  G3  size zip zone zoom
  G4  size zone zoom zero
  G5  magazine magazines arizona brazil
  G6  size zip zone zoom zero

\end{Verbatim}
\section{Dvorak Corpus by Position Matching}
\begin{Verbatim}
Exercise 1   new keys: aehnostu
  G1  aoeu snth
  G2  aoeua snths
  G3  the that not has
  G4  that these than state
  G5  states estate status season
  G6  the to on that not

Exercise 2   new keys: di
  G1  aoeui snthd
  G2  aoeuiua snthdhs
  G3  and this site his
  G4  this site date into
  G5  united students issues student
  G6  and in is this it

Exercise 3   new keys: .g
  G1  e.e hgh
  G2  e.e hGh
  G3  get good high sign
  G4  good high using sign
  G5  design things against hosting
  G6  get go good high sign

Exercise 4   new keys: pr
  G1  upu nrn
  G2  uPu nRn
  G3  are page our here
  G4  page other their there
  G5  support rights through report
  G6  or are page our up

Exercise 5   new keys: fy
  G1  uyu hfh
  G2  uYu hFh
  G3  for you your free
  G4  your free they first
  G5  history property different features
  G6  of for you your if

Exercise 6   new keys: ,c
  G1  o,o tct
  G2  o,o tCt
  G3  can city such each
  G4  price city such each
  G5  search contact products product
  G6  can city such each car

Exercise 7   new keys: bm
  G1  hbh hmh
  G2  hBh hMh
  G3  from more home but
  G4  from more home about
  G5  business system number message
  G6  by be from more home

Exercise 8   new keys: kx
  G1  uku uxu
  G2  uKu uXu
  G3  back next make book
  G4  back next make books
  G5  market feedback making marketing
  G6  back next make book take

Exercise 9   new keys: jw
  G1  eje twt
  G2  eJe tWt
  G3  with new was what
  G4  with what which news
  G5  software within subject between
  G6  with new was we what

Exercise 10   new keys: 'l
  G1  a'a sls
  G2  a'a sLs
  G3  all will only also
  G4  will only also help
  G5  online people health should
  G6  all will only also help

Exercise 11   new keys: qv
  G1  oqo nvn
  G2  oQo nVn
  G3  have view over very
  G4  have view over video
  G5  services service available review
  G6  have view over very even

Exercise 12   new keys: z
  G1  szs
  G2  sZs
  G3  size zip zone zoom
  G4  size zone zoom zero
  G5  magazine magazines arizona brazil
  G6  size zip zone zoom zero

Exercise 13   new keys: ;
  G1  a;a
  G2  a;a
  G3  (no listed word contains a new key)
  G4  (no listed word contains a new key)
  G5  (no listed word contains a new key)
  G6  (none)

\end{Verbatim}
\section{Letter Matching}

Under letter matching the word groups G3 to G6 of the QWERTY corpus are reused unchanged, since $\Vocab_i$ is the same. Only the chains differ, because the sandwich drills must be rebuilt from the home keys of the Dvorak layout for the physical positions that the letters occupy there.

